\documentclass[UTF8,11pt,a4paper,fontset=fandol]{ctexart}
\usepackage[margin=20mm]{geometry}
\usepackage{amsmath,amssymb,booktabs,array,tabularx,enumitem,setspace,xcolor,hyperref,xurl,fancyhdr}
\setstretch{1.08}
\setlength{\emergencystretch}{3em}
\setlength{\headheight}{15pt}
\setlength{\tabcolsep}{5pt}
\renewcommand{\arraystretch}{1.2}
\setlist{nosep,leftmargin=2em}
\hypersetup{colorlinks=true,linkcolor=blue!45!black,urlcolor=blue!45!black}
\providecommand{\WorkloadName}{工作负载}
\pagestyle{fancy}\fancyhf{}
\fancyhead[L]{\small Table II(b)：算子级 RI}
\fancyhead[R]{\small\WorkloadName}
\fancyfoot[C]{\thepage}
\newcommand{\qs}{Q_{\mathrm S}}
\newcommand{\qr}{Q_{\mathrm R}}
\newcommand{\RI}{\mathrm{RI}}
\newcommand{\para}[1]{\par\smallskip\noindent\textbf{#1}\quad}
\newcommand{\ModelTableSetup}{\renewcommand{\arraystretch}{1.2}\setlength{\tabcolsep}{4pt}}
\newcommand{\ResultTableSetup}{\renewcommand{\arraystretch}{1.45}\setlength{\tabcolsep}{4pt}}
\newcommand{\DisplayNote}{\par\smallskip{\noindent\footnotesize $1\mathrm{K}=1024$，$1\mathrm{M}=1024^2$。RI 至少为 1 时保留一位小数，小于 1 时保留三位有效数字；精确值见机器数据。\par}\smallskip}
\author{}\date{2026 年 10 月 2 日}

% Compact, identical title block keeps the mathematical introduction on one page.
\makeatletter
\renewcommand{\maketitle}{\thispagestyle{plain}\begin{center}
{\LARGE\bfseries\@title\par}\vspace{8pt}{\normalsize\@date\par}
\end{center}\vspace{3pt}}
\makeatother

\renewcommand{\WorkloadName}{QKV 投影}
\title{\bfseries\WorkloadName 的输入复用与 RI}
\begin{document}
\maketitle

\section{对象与公式}
取每个模型的一层完整／全局 GQA。将 Q、可选的输出门控 G、K、V
投影权重完整装载一次，累计服务 $U$ 个输入向量。沿用 II(a) 的复用次数（reuse count）：包含首次使用，可跨多个 batch 或请求累计。
输入与权重均按每元素 1 Byte 计数。各投影使用同一个 $X\in\mathbb R^{U\times D}$，
权重形状为 $W_j\in\mathbb R^{N_j\times D}$，输出为 $XW_j^{\mathsf T}$。
本对象到这组投影结束，不包含后续 O 投影。

原生 query 与 KV 头数分别为 $H_q,H_{kv}$，每头 QK 与 V 维度为 $d_{QK},d_V$；
额外门控宽度为 $D_g$，未启用时取 0。于是
\[
\begin{aligned}
N_Q&=H_qd_{QK}, & N_G&=D_g,\\
N_K&=H_{kv}d_{QK}, & N_V&=H_{kv}d_V,\\[-2pt]
N_{\rm proj}&=N_Q+N_G+N_K+N_V.&&
\end{aligned}
\]
同一矩阵阶段的共享输入只计一次；各权重写入相加。沿用 Table II(a) 的一次装载口径，得到
\[
\boxed{\displaystyle
\qs=UD,\qquad \qr=DN_{\rm proj},\qquad
\RI=\frac{\qs}{\qr}=\frac{U}{N_{\rm proj}}.}
\]
若共有 $m$ 个投影分项，逐项独立记录时各有 $\qs^{(j)}=UD$、$\qr^{(j)}=DN_j$；
汇总输入需扣除 $(m-1)UD$ 的同源重叠。投影输出不直接加入本阶段的 $\qs$。

\section{模型信息}
下表列出固定层的原生投影宽度；层号从 0 起，所有矩阵的输入维度均为 $D$。
$N_G=0$ 表示没有额外 G 投影。
\begin{center}\small\ModelTableSetup
\begin{tabularx}{\linewidth}{@{}p{0.32\linewidth}*{7}{>{\centering\arraybackslash}X}@{}}
\toprule
模型 & 层号 & $D$ & $N_Q$ & $N_G$ & $N_K$ & $N_V$ & $N_{\rm proj}$ \\
\midrule
Qwen3.5-2B & 3 & 2048 & 2048 & 2048 & 512 & 512 & 5120 \\
Ministral 3 8B (2512) & 0 & 4096 & 4096 & 0 & 1024 & 1024 & 6144 \\
Qwen3.6-35B-A3B & 3 & 2048 & 4096 & 4096 & 512 & 512 & 9216 \\
Tencent Hy3 (295B) & 1 & 4096 & 8192 & 0 & 1024 & 1024 & 10240 \\
Ling-1T & 4 & 8192 & 8192 & 0 & 1024 & 1024 & 10240 \\
MiMo-V2.5-Pro & 7 & 6144 & 24576 & 0 & 1536 & 1024 & 27136 \\
\bottomrule
\end{tabularx}
\end{center}
两款 Qwen 的原生 \texttt{q\_proj} 同时产生 Q 与 G，表中分别保留其逻辑宽度。
MiMo 的 $d_{QK}=192$、$d_V=128$，故 K 与 V 的宽度不同；其余五款的 QK 与 V 每头维度相同。
固定版本与原始实现的定位保存在机器配置中。

\newpage
\section{代入结果}
表内为小数 $\RI$；前四列为 $U\in\{1,1\mathrm{K},128\mathrm{K},1\mathrm{M}\}$，末列为无界复用极限。
\begin{center}\small\ResultTableSetup
\begin{tabularx}{\linewidth}{@{}p{0.32\linewidth}*{5}{>{\centering\arraybackslash}X}@{}}
\toprule
模型 & $U=1$ & $U=1\mathrm{K}$ & $U=128\mathrm{K}$ & $U=1\mathrm{M}$ & $U\to\infty$ \\
\midrule
Qwen3.5-2B & $0.000195$ & $0.200$ & $25.6$ & $204.8$ & $\infty$ \\
Ministral 3 8B (2512) & $0.000163$ & $0.167$ & $21.3$ & $170.7$ & $\infty$ \\
Qwen3.6-35B-A3B & $0.000109$ & $0.111$ & $14.2$ & $113.8$ & $\infty$ \\
Tencent Hy3 (295B) & $0.0000977$ & $0.100$ & $12.8$ & $102.4$ & $\infty$ \\
Ling-1T & $0.0000977$ & $0.100$ & $12.8$ & $102.4$ & $\infty$ \\
MiMo-V2.5-Pro & $0.0000369$ & $0.0377$ & $4.8$ & $38.6$ & $\infty$ \\
\bottomrule
\end{tabularx}
\end{center}
\DisplayNote
例如 Qwen3.5-2B 在 $U=1\mathrm{K}$ 时，
\[
\begin{aligned}
N_{\rm proj}&=2048+2048+512+512=5120,\\
\qs&=1024\times2048=2\,097\,152\ \mathrm{Byte},\\
\qr&=2048\times5120=10\,485\,760\ \mathrm{Byte},\\
\RI&=\frac{2\,097\,152}{10\,485\,760}=0.200.
\end{aligned}
\]

\section{结果说明}
固定模型时，$U$ 每增加一倍，$\qs$ 与 $\RI$ 都增加一倍，而一次权重写入量 $\qr$ 不变。
$U=1$ 是这份装载只服务一个向量的端点；$U\to\infty$ 时，输入量与 RI 无界增长，
但 $\qr$ 始终为有限的 $DN_{\rm proj}$，有效权重容量也保持该值。

跨模型比较由所选投影的总输出宽度 $N_{\rm proj}$ 决定。
Hy3 与 Ling 的 $N_{\rm proj}$ 都为 10240，因而各档 RI 相同；
Ling 的 $D$ 是 Hy3 的两倍，绝对输入与权重字节也都翻倍，比例恰好抵消。
模型的总参数标签不能替代这组投影尺寸。

额外 G 投影增加权重写入量，却继续共享同一份输入，因此必须计入 $N_{\rm proj}$。
MiMo 的原生 Q/K/V 总宽度为 27136，是本组六款中最大的，故相同 $U$ 下 RI 最小；
这一差别来自实际头数及头维，而非模型名称或输入维度 $D$ 本身。

\par\smallskip\noindent\small
固定模型来源、精确字节及复算记录见本目录 README 与 data/；数学入口为 Table II(a)。
\end{document}
